\section{Subsets}
It is common that every element, or some elements of a set $A$ are also elements of another set $B$. Consider the following: 
$$
A=\{ 1,2,3 \}, B=\{ 0,1,2,3,4,5 \}.
$$
In cases as this one, we say that $A\subseteq B$, we might even say that $A\subset B$ (since all elements of $A$ are contained within $B$).

\begin{defn}
    Suppose $A,B$ are sets. If every element of $A$ is also an element of $B$, then we say that $A\subseteq B$. We write $A\not\subseteq B$ if $A$ is \emph{not} a subset of $B$ i.e., if it is not true that \emph{every} element of $A$ is also an element of $B$. So $A\not\subseteq B$ tells us that \emph{at least one} element of the set on the left is not contained in the set on the right.
\end{defn}

\begin{example}
    Some key examples are the following: 
    \begin{enumerate}
        \item $\{ 2,3,7 \}\subseteq \{ 2,3,4,5,6,7 \}$;
        \item $\{ 2,3,7 \}\not\subseteq \{ 2,3,4,5,6 \}$;
        \item $\N\subseteq\Z\subseteq\Q\subseteq\R$;
        \item $\R\times\N\subseteq\R^2$;
        \item $A\subseteq A$, any set $A$ is a subset of itself; 
        \item $\varnothing\subseteq\varnothing$.
    \end{enumerate}
\end{example}

These examples bring out an important fact about the empty set. 
\begin{center}
 \boxed{\text{If } B \text{ is any set whatsoever, then } \varnothing\subseteq B}.   
\end{center}

\begin{fact}
    The empty set is a subset of all sets, that is, $\varnothing\subseteq B$ for any set $B$.
    \begin{proof}
        Let $B$ be any set, then consider the definition of subset. Then if $\varnothing\not\subseteq B$ this would mean that there are \emph{some} members of $\varnothing$ which are not members of $B$. However, we know that by definition $x\notin\varnothing$, or more formally 
        $$
        \varnothing = \{ x:x\ne x \}.
        $$
        Therefore, $\varnothing\subseteq B$.
    \end{proof}
    \emph{Comment:} This proof technique is called \emph{Proof by Contradiction} and it does come up a lot in the exercises later on, so it is worth trying to understand what is happening at this stage, however, the proof strategy will be explained in detail later. 
\end{fact}

\begin{fact}
    If $\lvert B \rvert=n$, then $B$ has $2^n$ subsets. Further, for \emph{any finite} set $B$ that has $n$ number of elements, then $B$ has $2^n$ subsets.    
\end{fact}

\begin{example}
    Consider the subsets of the set $B=\{ 1,2,\{ 1,3 \} \}$, there are $3$ elements so $2^3$. 
    $$
    \{\varnothing,\{ 1 \},\{ 2 \}, \{ \{ 1,3 \} \}, \{ 1,2 \},\{ 1,2,\{ 1,3 \} \}, \{ 1,\{ 1,3 \} \},\{ 2,\{ 1,3 \} \}\}.
    $$
    
    \emph{Important:} It is really important to notice that $\{ 1,3 \}$ is \textbf{not} a subset of $B$. As in: 
    $$
    \{ 1,3 \}\not\subseteq B \text{ but }\{ 1,3 \}\in B, \text{ however, } \{ \{ 1,3 \} \}\subseteq B.
    $$
\end{example}

\subsection{Exercises from Section 3}
\begin{exercise}\label{ex:3.1}
    List all the subsets of the following sets:
    \begin{enumerate}
        \item $\{ 1,2,3,4 \}$
        \item $\{ 1,2,\varnothing \}$
        \item $\{ \{ \R \} \}$
        \item $\varnothing$
        \item $\{ \varnothing \}$
        \item $\{ \R,\N,\Q \}$
        \item $\{ \R,\{ \N,\Q \} \}$
        \item $\{ \{ 0,1 \},\{ 0,1,\{ 2 \} \},\{ 0 \} \}$
    \end{enumerate}
\end{exercise}\hyperref[sol:3.1]{[Go to Solution]}


\begin{exercise}\label{ex:3.2}
    Write out the following sets by listing their elements between braces.
    \begin{enumerate}
        \item $\{ X:X\subseteq \{ 3,2,a \}\text{ and } \lvert X \rvert=2 \}$
        \item $\{ X\subseteq\N:\lvert X \rvert\le 1 \}$
        \item $\{ X:X\subseteq \{ 3,2,a \}\text{ and }\lvert X \rvert=4 \}$
        \item $\{ X:X\subseteq \{ 3,2,a \}\text{ and }\lvert X \rvert=1 \}$
    \end{enumerate}
\end{exercise}\hyperref[sol:3.2]{[Go to Solution]}


\begin{exercise}\label{ex:3.3}
    Decide if the following statements are true or false. Explain. 
    \begin{enumerate}
        \item $\R^3\subseteq\R^3$
        \item $\R^2\subseteq\R^3$
        \item $\{ (x,y)\in\R^2:x-1=0 \}\subseteq \{ (x,y)\in\R^2:x^2-x=0 \}$
        \item $\{ (x,y)\in\R^2:x^2-x=0 \}\subseteq \{ (x,y)\in\R^2:x-1=0 \}$
    \end{enumerate}
\end{exercise}\hyperref[sol:3.3]{[Go to Solution]}